\documentclass[conference]{IEEEtran}
\IEEEoverridecommandlockouts
% The preceding line is only needed to identify funding in the first footnote. If that is unneeded, please comment it out.
%
% Class Project Report -- EN.665.715 Embodied Agentic AI for Robotic Control
% Built from the instructor's "Class Project Template" (conference_101719.tex, IEEEtran).
% DRAFT STATUS: only Section II (Related Works) and the References are written.
% AI assistance: Claude Fable 5.1 via Claude Code drafted this file; see the disclosure
% section before the references.
\usepackage{cite}
\usepackage{amsmath,amssymb,amsfonts}
\usepackage{algorithmic}
\usepackage{graphicx}
\usepackage{textcomp}
\usepackage{xcolor}
\usepackage{comment}
% caption=false keeps IEEEtran's own caption format (centered small-caps table titles);
% the template loaded subfig without it, which lets the caption package override that format.
\usepackage[caption=false,font=footnotesize]{subfig}
% Balances the two columns of the last page (the reference list).
\usepackage{flushend}
\newcommand{\centered}[1]{\begin{tabular}{l} #1 \end{tabular}}
\def\BibTeX{{\rm B\kern-.05em{\sc i\kern-.025em b}\kern-.08em
    T\kern-.1667em\lower.7ex\hbox{E}\kern-.125emX}}
% Marker for the sections that this draft does not cover.
\newcommand{\laterdraft}{\textit{Reserved for a later draft.}}
\begin{document}
% IEEE reference style: more than six authors are printed as "first author et al."
\bstctlcite{IEEEexample:BSTcontrol}

\title{BeaconMind: Can a Large Language Model Own Landmark SLAM?\\
%{\footnotesize \textsuperscript{*}Note: Sub-titles are not captured %in Xplore and
%should not be used}
%\thanks{Identify applicable funding agency here. If none, delete this.}
}

\author{\IEEEauthorblockN{Xiaohui Chen}
\IEEEauthorblockA{\textit{Engineering for Professionals, Whiting School of Engineering} \\
\textit{Johns Hopkins University}\\
Baltimore, United States \\
EN.665.715 Embodied Agentic AI for Robotic Control (individual project)}
}

\maketitle

\begin{abstract}
Reserved for a later draft. This submission covers only Section~\ref{sec:related}
(Related Works) and the References; the other sections are placeholders.
\end{abstract}

\begin{IEEEkeywords}
Reserved for a later draft.
\end{IEEEkeywords}

\section{Introduction}
\laterdraft

\section{Related Works}\label{sec:related}
BeaconMind, the system proposed in this project, asks which parts of landmark-based
simultaneous localization and mapping (SLAM) a large language model (LLM) can own, that
is, for which parts the model itself can make the estimation decisions. A low-cost
differential-drive car with wheel encoders, an inertial measurement unit, and one
monocular camera, but no range sensor, explores a walled arena whose only landmarks are
AprilTags printed on paper and fixed to the walls. Each physical tag is a \emph{beacon};
because tag identifiers may repeat in the benchmark, an identifier does not always name one
beacon. The car moves in a stop-observe-decide-move cycle. At each stop, an LLM agent
receives the odometry increment and the tag detections as text, obtains every geometric
quantity from deterministic tools, and maintains its estimate of the car's pose (the pose
belief) and the beacon map: it decides which detections to trust, which beacon a detection
belongs to, whether a pose fix or a loop closure (a recognized revisit) is accepted, and
which of a few fixed, short moves to make next. The project is at the design stage, so
statements about BeaconMind describe the planned system and its planned evaluation, not
results. Four strands of prior work shape the design; for each, this section summarizes
what was shown, what it leaves open, and how BeaconMind differs.

\subsection{Language-Conditioned Planning and Closed-Loop Agents}
A first line of work made LLM plans executable by grounding them in what the robot can
do. Huang \textit{et~al.} showed that a pretrained LLM decomposes a household task into
plausible steps, but that naively generated plans often cannot be mapped to admissible
actions~\cite{huang2022zeroshot}. SayCan addressed this by weighting each candidate skill
with a learned affordance value; removing that grounding lowered planning success from
84\% to 67\% over 101 instructions~\cite{ahn2022saycan}. Inner Monologue then returned
success detection and scene descriptions to the planner as text and found that
closed-loop language feedback significantly improves instruction
completion~\cite{huang2023inner}. ReAct formulated the pattern in general form,
interleaving reasoning, action, and observation, and outperformed imitation and
reinforcement learning baselines on the ALFWorld and WebShop benchmarks by 34 and 10
percentage points in success rate, respectively~\cite{yao2023react}. Kambhampati
\textit{et~al.} argue that this line of work has a limit: LLMs can neither plan nor verify
plans by themselves and, in what the authors call an LLM-Modulo framework, should generate
candidates that external, sound verifiers check~\cite{kambhampati2024llmmodulo}.

These ideas are now packaged for robot software. ROSA wraps the Robot Operating System
(ROS) in a ReAct agent with predefined tools so that operators can command robots in
natural language~\cite{royce2025rosa}, and the survey by Salimpour \textit{et~al.}
classifies such integrations from single-shot protocol translation, through ReAct-style
agentic interfaces and orchestration, to end-to-end embedded
policies~\cite{salimpour2025towards}. Benchmarks quantify what still fails. The Embodied
Agent Interface reports that trajectory-feasibility errors are common (45.2\%), a large
portion of them missing-step and additional-step errors~\cite{li2024embodied}, and in
SafeAgentBench the most safety-conscious agent rejects only 10\% of detailed hazardous
tasks~\cite{yin2025safeagentbench}.

\begin{table*}[t]
\caption{Division of Labor Between the Language Model and Classical Geometry in Representative Prior Systems and in BeaconMind}
\label{tab:division}
\centering
\footnotesize
\renewcommand{\arraystretch}{1.1}
\begin{tabular}{p{0.17\textwidth} p{0.36\textwidth} p{0.40\textwidth}}
\hline
\textbf{System} & \textbf{Pose and map supplied by} & \textbf{What the language model contributes} \\
\hline
SayPlan~\cite{rana2023sayplan} & Pre-constructed 3D scene graph & Task planning over the scene graph \\
VLFM~\cite{yokoyama2024vlfm} & Odometry sensor; occupancy map built from depth & Scores which frontier to explore (vision-language model) \\
L3MVN~\cite{yu2023l3mvn} & Odometry sensor; semantic map from RGB-D images & Selects the frontier relevant to the target object \\
Lee \textit{et~al.}~\cite{lee2026spatial} & Cartographer SLAM (Nav2 executes the goals) & Picks goals from annotated map images; logs semantics \\
Syarubany \textit{et~al.}~\cite{syarubany2026llm} & ORB-SLAM3 global pose; AprilTags as anchors & One constrained discrete command at each junction \\
SEO-SLAM~\cite{hong2025seoslam} & Odometry, linear-assignment association, factor graph & Object verification, label refinement, and landmark merging \\
Sandfuchs \textit{et~al.}~\cite{sandfuchs2026llms} & Oracle localization in a text gridworld & One grid move per step, to explore or to navigate \\
\hline
BeaconMind (this project, planned) & The LLM's pose belief and beacon map, updated at every stop with single-timestep geometry tools; in the primary design a classical factor graph serves only as the reference & Which detections to trust, data association under duplicated tag identifiers, acceptance of pose fixes and loop closures, map updates, and the choice of each move \\
\hline
\end{tabular}
\end{table*}

BeaconMind adopts these lessons. Its agent is a ReAct-style agentic interface in the
survey's classification, and ROSA is its planned demonstration runtime; the benchmark will
instead use a minimal harness that exposes only the project's own tools, because stock
ROSA adds its own tools and fixed prompts to every call and would confound a comparison
of models. Its actions are admissible by construction, because the agent may choose only
from a small fixed set of short moves, and its loop is closed, because every decision
starts from a fresh observation. The missing-step and additional-step errors that the
Embodied Agent Interface reports cannot arise in that form, since the agent emits one
move per observation instead of a multi-step plan. Following the LLM-Modulo argument, a
deterministic validator between the model and the motors rejects any move outside that set
or beyond the safety limits. Because the agent also accepts free-text operator requests,
its refusal of unsafe ones is proposed to be measured with a fixed request suite and a
rejection rate adapted from SafeAgentBench. The difference lies in what the model reasons
about: the works above plan \emph{tasks} and take the robot's state as given, whereas
BeaconMind asks the model to estimate that state, the pose and the map that such plans
presuppose.

\subsection{Language Models in Navigation, Mapping, and SLAM}
In the language-driven navigation systems reviewed here, the geometry stays classical
(Table~\ref{tab:division}). SayPlan plans over a pre-constructed 3D scene graph and
delegates path planning to Dijkstra's algorithm; its authors report that the LLM fails at
simple distance-based reasoning and name the pre-built, static graph as a
limitation~\cite{rana2023sayplan}, whereas BeaconMind's map does not exist before the run
and is built by the agent. For unknown spaces, VLFM builds an occupancy map from depth
observations and uses a vision-language model to score frontiers (the boundaries between
explored and unexplored space) and pick the next one to visit~\cite{yokoyama2024vlfm},
and L3MVN has an LLM select the relevant frontier from a semantic map built from RGB-D
images~\cite{yu2023l3mvn}; both report strong object-goal navigation results in
simulation and a real-robot demonstration. Both also need depth to build the map and take
the robot's position from an odometry sensor; BeaconMind has no range sensor, so its map
is a set of beacon coordinates, not an occupancy grid. Lee \textit{et~al.} connect LLMs to
ROS through Model Context Protocol tools that render the occupancy grid as annotated
images; with Cartographer for SLAM and Nav2 for navigation in Isaac Sim, Claude Sonnet 4.6
and GPT-5.5 reach 97.8\% and 99.8\% map coverage and succeed in 26 and 28 of 30
spatial-reasoning navigation trials, respectively~\cite{lee2026spatial}. The closest
system identified in the literature search is that of Syarubany \textit{et~al.}: an LLM
issues one constrained discrete command at each corridor junction, and AprilTags serve as
stable anchors, but ORB-SLAM3 provides the global pose and the evaluation of the
integrated system is qualitative~\cite{syarubany2026llm}.

Where language models come closer to SLAM itself, they supply semantics or discrete
choices, not estimates. SEO-SLAM assigns multimodal LLM agents to object verification,
label refinement, and landmark merging, and reports higher semantic accuracy and fewer
false positives than its baselines, while data association (deciding which landmark an
observation belongs to) is solved as a linear assignment problem over semantic and
Mahalanobis distances and the map comes from factor-graph optimization, the standard
graph-based formulation of SLAM~\cite{hong2025seoslam}. Lyu \textit{et~al.} let a
multimodal LLM make the final choice among place-recognition candidates retrieved by a
vision foundation model~\cite{lyu2024tell}; this is a language-based form of the place
recognition behind loop closure, but its output is a choice among images, not a pose.
Sandfuchs \textit{et~al.} evaluate nine LLMs as text-only controllers in partially
observable gridworlds in which the agent's coordinates are always known, and they state
the boundary explicitly: ``This is not SLAM: the uncertainty is in the map, not in
pose''~\cite{sandfuchs2026llms}. The strongest reasoning-tuned models (GPT-5 and
GPT-OSS-120B) reach the goal in all three layouts, but with the pose supplied, which is
the condition BeaconMind removes. Learned visual navigation policies such as NoMaD
sidestep metric estimation altogether: they use no language model, turn camera images
directly into actions and, for long horizons, rely on a topological graph of past
observations instead of a metric map~\cite{sridhar2024nomad}. BeaconMind lists such a
policy as a candidate exploration baseline; it cannot be scored on pose or map error,
because it yields no metric pose or map.

In none of the language-model systems above is the robot's pose estimate in the model's
hands. SEO-SLAM comes closest on the map side: its agents prune and merge object landmarks
produced by a conventional detection, association, and factor-graph pipeline and feed
those edits back into the graph. The paper scores the semantic precision and recall of
the resulting map and the absolute pose error of the optimized trajectory, but the
estimate remains the optimizer's. None of these systems can say how accurate a pose or a
geometric map maintained by a language model would be. BeaconMind moves that responsibility to the LLM: deterministic tools still
do the arithmetic, but the LLM decides at every stop which results enter the pose belief
and the beacon map. In the primary design the classical estimator is kept only as a
reference (variants with no tools, or with the LLM on top of the classical estimator, are
proposed as stretch conditions), and the LLM's belief is scored against ground truth at every stop. To
my knowledge, no prior study reports a trajectory or map error for a landmark-SLAM
estimate maintained by an LLM. That measurement is the gap this project targets, and a
negative outcome would still be informative.

\subsection{Limits of LLM Spatial Reasoning and the Case for Tools}
Recent diagnostics show why handing estimation to a language model is risky, and they
shape how BeaconMind does it. Jiang \textit{et~al.} find that single-call, whole-route
maze navigation by GPT-4o, DeepSeek-V3, and Llama-3.3-70B falls to near-zero success by
$10\times10$ mazes; that structured coordinate text is the most navigable input format,
far ahead of rendered images; and that handing corridor walking and dead-end retreat to a
deterministic walker, with the model queried only at junctions through an explicit
cell-type prompt, lifts GPT-4o's success by up to 92 percentage points at mid sizes,
although the collapse returns by $30\times30$~\cite{jiang2026lost}. Li \textit{et~al.}
find that models treat each spatial question independently instead of building cumulative
spatial knowledge, and that accuracy depends on the representation: Gemini-2.5-Flash with
chain-of-thought prompting reaches \mbox{80--86\%} on small mazes given as adjacency lists
but \mbox{16--34\%} when they are given as visual character grids~\cite{li2026llms}. For
perception, Deng \textit{et~al.} show that the best vision-language model reaches a
macro-F1 of 0.66 on relative camera pose estimation, where a geometric pipeline (LoFTR)
reaches 0.99 and humans 0.91~\cite{deng2026lost}. These results do not contradict Lee
\textit{et~al.}, whose models read annotated map images to pick goal positions and judge
coverage while a classical planner executed every path. On the constructive side, PAL
shows that LLMs decompose problems well but make arithmetic mistakes even when the
decomposition is correct, and that off-loading the solution step to a program interpreter
gives more accurate results than much larger models~\cite{gao2023pal}.

BeaconMind is designed around these findings. The model receives the tag detector's output
(identifier, range, bearing, and a quality score) as structured text rather than pixels,
with raw-image input kept as a separate planned ablation, and it keeps its state in
explicit coordinates. As in PAL, every geometric quantity comes from a deterministic tool;
unlike PAL's interpreter, which solves the whole problem, each tool is restricted to a
single timestep: it is a function of one stop's detections, the previous belief, and a
few harness counters, for example composing one odometry increment, converting one range
and bearing into a map point, or fixing the pose from the beacons seen at that stop. None
optimizes over the observation history. The system prompt prescribes a per-step procedure
(predict, associate, correct the pose, update the map, check loop closure, choose a move),
and the model carries it out by calling the tools and acting on their results, so
trusting, matching, and accepting remain steps that the model must execute correctly.
Because that procedure can also be run as ordinary code, a rule-program baseline is
proposed to measure what executing it with an LLM costs in accuracy, time, and money; what
the LLM adds beyond the procedure, such as handling free-text requests and refusing or
asking, would be scored by the request suite. Each
decision is a fresh, stateless conversation that carries only the mission, any operator
request, the model's previous report, and the new observation. Every step is thus posed as an isolated
decision, which Jiang \textit{et~al.} measured to be more reliable than the same decision
made inside a long plan. Whether that is enough is open: in their study the collapse
returned even with per-junction queries, and the maze studies used symbolic grids with
perfect observations, whereas BeaconMind poses a physical estimation problem with noisy
odometry and noisy range and bearing measurements.

\subsection{Fiducial Landmark SLAM and Evaluation Practice}
Fiducial SLAM supplies both the landmarks and the reference solution. AprilTag is a visual
fiducial system that gives a tag's identity and its full 6-DOF pose relative to the camera
from a single image~\cite{olson2011apriltag}. TagSLAM builds a factor graph over tags,
cameras, and bodies on the GTSAM factor-graph library and observes that, with a small set
of artificial landmarks, ``loop closure is guaranteed once a tag is successfully
decoded''~\cite{pfrommer2019tagslam}; iSAM2 solves such graphs incrementally through the
Bayes tree~\cite{kaess2012isam2}. For evaluation, Sturm \textit{et~al.} formalize the
absolute trajectory error (ATE) and the relative pose error (RPE) as the trajectory
metrics of their SLAM benchmark~\cite{sturm2012benchmark}, and Anderson \textit{et~al.}
recommend that a navigation agent signal completion with an explicit ``DONE'' action and
be scored by success weighted by path length (SPL)~\cite{anderson2018evaluation}.

BeaconMind uses this body of work in two ways. First, fiducial SLAM supplies the
yardstick: a GTSAM/iSAM2 planar factor graph is the best-case reference and dead reckoning
(odometry alone) is the floor, and the LLM's ATE, RPE, and beacon-map error will be
reported against both on identical replayed logs from simulation and from the real car. In
replay the moves are those recorded in the log and the model only estimates; move choice
is evaluated separately in closed loop, where navigation with the finished map will be
scored by ground-truth success at ``DONE'' and by SPL. Second, BeaconMind removes the
assumption that makes data association in fiducial SLAM trivial. TagSLAM requires unique
identifiers, so decoding a tag settles data association and, with it, loop closure.
BeaconMind's benchmark adds a condition in which identifiers are duplicated on purpose, so
that association and, when a revisit occurs, loop closure require reasoning over geometry
and the map built so far. For that condition it is proposed that the reference receive the
true association, which makes it an upper bound on achievable accuracy, not a like-for-like
competitor. The systems in Table~\ref{tab:division} either face no landmark association or
solve it in classical code. SEO-SLAM is a partial exception: its measurement-to-landmark
association is a linear assignment, but its language-model agents afterwards merge
duplicate landmarks and remove erroneous ones. Whether an LLM can make these decisions
with useful precision, and at what latency and cost, is the question BeaconMind is built
to answer.

\section{System Overview}\label{sytem_overview}
\laterdraft

\section{Experiments}\label{experiments}
\laterdraft

\section{Results and Evaluation}\label{results}
\laterdraft

\section{Conclusion}
\laterdraft

\section*{Generative AI Disclosure}
Claude Fable 5.1 (Anthropic), used through Claude Code, assisted with this draft: it
searched for and retrieved the cited papers, checked each summarized claim against the
paper's text, drafted the wording of Section~\ref{sec:related} and
Table~\ref{tab:division}, prepared the \BibTeX{} entries and the \LaTeX{} source, and ran
an adversarial review of the draft. Many of the sources were first identified in the
project's earlier AI-assisted literature searches (Claude Opus 5.5 in Claude Code and in
Cursor, with citations re-checked by Grok 4.6). The project's design documents, on which
the descriptions of BeaconMind rest, were drafted, and many of their detailed design
choices made, by Claude Opus 5.5 under the author's delegation; the author's review of
some of those choices is still in progress. The author chose the project direction,
reviews all AI-assisted content against the sources before submission, and is responsible
for its accuracy and originality.

%\section*{Acknowledgment}

\bibliographystyle{IEEEtran}
\bibliography{IEEEabrv,mybibfile}

\end{document}
